% from ega4-2-fr.tex
\oldpage[IV-2]{226}
\phantomsection
\addcontentsline{toc}{section}{Index terminologique}
\section*{INDEX TERMINOLOGIQUE}

\begingroup
\footnotesize
\setlength{\parindent}{0pt}
\setlength{\parskip}{1.1pt}
\newcommand{\EGAIVtwoTerm}[2]{#1 : #2.\par}

\EGAIVtwoTerm{Algèbre séparable sur un corps}{4.6.2}
\EGAIVtwoTerm{Anneau des séries formelles restreintes sur un corps valué complet}{7.4.8}
\EGAIVtwoTerm{Anneau excellent}{7.8.2}
\EGAIVtwoTerm{Anneau formellement caténaire}{7.1.9}
\EGAIVtwoTerm{Anneau formellement équidimensionnel}{7.1.1}
\EGAIVtwoTerm{Anneau géométriquement normal, géométriquement réduit,
géométriquement régulier, possédant la propriété $(R_k)$ géométrique}{6.7.6}
\EGAIVtwoTerm{Anneau normal}{5.13.5}
\EGAIVtwoTerm{$P$-anneau}{7.3.13}
\EGAIVtwoTerm{Anneau possédant la propriété $(R_k)$}{5.8.2}
\EGAIVtwoTerm{Anneau strictement équidimensionnel}{7.2.1}
\EGAIVtwoTerm{Anneau strictement formellement caténaire}{7.2.6}
\EGAIVtwoTerm{Anneau universellement caténaire}{5.6.2}
\EGAIVtwoTerm{Antifiltre de parties d'un ensemble}{5.10.8}

\medskip
\EGAIVtwoTerm{Codimension d'une partie quelconque d'un préschéma}{5.1.3}
\EGAIVtwoTerm{Coprofondeur d'un Module en un point, coprofondeur d'un Module,
coprofondeur d'un préschéma}{5.7.1}
\EGAIVtwoTerm{Corps de définition}{4.8.4}
\EGAIVtwoTerm{Corps séparablement clos}{4.3.3}
\EGAIVtwoTerm{Cycle premier associé à un Module, à un préschéma}{3.1.1}
\EGAIVtwoTerm{Cycle premier associé immergé}{3.1.1}

\medskip
\EGAIVtwoTerm{Décomposition irredondante, décomposition irredondante réduite
d'un Module}{3.2.5}
\EGAIVtwoTerm{Décomposition primaire de $0$ dans un Module}{3.2.5}
\EGAIVtwoTerm{Dimension d'un Module}{5.1.12}
\EGAIVtwoTerm{Dimension d'un préschéma sur un corps}{4.1.1}

\medskip
\EGAIVtwoTerm{Exposant d'inséparabilité d'une extension d'un corps}{4.7.4}
\EGAIVtwoTerm{Extension primaire d'un corps}{4.3.1}

\medskip
\EGAIVtwoTerm{Fermeture algébrique séparable d'un corps dans une extension}{4.3.4}
\EGAIVtwoTerm{Fibre formelle d'un anneau en un point de son spectre}{7.3.13}

\medskip
\EGAIVtwoTerm{Homéomorphisme universel}{2.4.2}
\EGAIVtwoTerm{Homomorphisme de Modules défini sur un corps}{4.8.4}

\medskip
\EGAIVtwoTerm{Lemme de Hironaka}{5.12.8}
\EGAIVtwoTerm{Longueur géométrique d'un Module en un point}{4.7.5}

\medskip
\EGAIVtwoTerm{Module intègre sur un anneau}{3.2.4}
\EGAIVtwoTerm{$A$-Module possédant la propriété $(S_k)$}{5.7.2}
\EGAIVtwoTerm{Module réduit sur un anneau}{3.2.2}
\EGAIVtwoTerm{$\mathcal O_X$-Module de Cohen-Macaulay en un point,
$\mathcal O_X$-Module de Cohen-Macaulay}{5.7.1}
\EGAIVtwoTerm{$\mathcal O_X$-Module fidèlement plat (relativement à un morphisme)}{2.2.4}
\EGAIVtwoTerm{$\mathcal O_X$-Module géométriquement intègre sur un corps}{4.6.19}

\clearpage
